\phantomsection
\section*{Глава 3. Реализация веб-приложения}
\addcontentsline{toc}{section}{Глава 3. Реализация веб-приложения}
\stepcounter{section}

\subsection{Инициализация приложения и развёртывание}

Приложение построено по паттерну Application Factory: функция
\texttt{create\_app()} создаёт экземпляр Flask, подключает расширения
SQLAlchemy, Flask-Migrate, Flask-Login и CSRFProtect, регистрирует пять
блупринтов и устанавливает контекстный процессор, передающий во все
шаблоны количество непрочитанных уведомлений текущего пользователя.

\begin{lstlisting}[caption={Функция Application Factory (app/\_\_init\_\_.py)}, label={lst:create_app}]
def create_app(config_name=None):
    if config_name is None:
        config_name = os.environ.get('FLASK_ENV', 'default')

    app = Flask(__name__)
    app.config.from_object(config[config_name])

    db.init_app(app);  migrate.init_app(app, db)
    login_manager.init_app(app);  csrf.init_app(app)

    from app.main        import main_bp
    from app.auth        import auth_bp
    from app.admin       import admin_bp
    from app.teacher     import teacher_bp
    from app.notifications import notifications_bp

    app.register_blueprint(main_bp)
    app.register_blueprint(auth_bp,           url_prefix='/auth')
    app.register_blueprint(admin_bp,          url_prefix='/admin')
    app.register_blueprint(teacher_bp,        url_prefix='/teacher')
    app.register_blueprint(notifications_bp,  url_prefix='/notifications')

    @app.context_processor
    def inject_notifications():
        if current_user.is_authenticated:
            count = Notification.query.filter_by(
                user_id=current_user.id, is_read=False
            ).count()
            return {'unread_count': count}
        return {'unread_count': 0}

    return app
\end{lstlisting}

Конфигурация передаётся через переменные окружения (\texttt{.env}): строка
подключения к PostgreSQL, секретный ключ сессии и корневой путь файлового
хранилища. Это позволяет запускать приложение в разных средах без
изменения исходного кода.

Среда выполнения описана в \texttt{docker-compose.yml}~\cite{docker_docs}
и включает два сервиса: \textbf{db} (PostgreSQL~16~\cite{postgresql_docs},
данные сохраняются в именованном томе) и \textbf{web} (Flask~3~\cite{flask_docs}
с Gunicorn, Python-код запекается в образ при сборке).
Миграции схемы управляются через Flask-Migrate (Alembic); команда
\texttt{flask db upgrade} выполняется при старте контейнера без пересборки образа.

\subsection{Аутентификация и ролевая модель}

Регистрация и вход реализованы в блупринте \texttt{auth}. Пароль
хешируется функцией \texttt{generate\_password\_hash} из модуля
\texttt{werkzeug.security} (алгоритм pbkdf2-sha256) и не хранится
в открытом виде. После входа Flask-Login
сохраняет идентификатор пользователя в подписанной серверной сессии
и предоставляет объект \texttt{current\_user} во всех маршрутах и шаблонах.

Контроль доступа реализован декораторами, определёнными в модуле
\texttt{app/auth/decorators.py} (листинг~\ref{lst:decorators}).

\begin{lstlisting}[caption={Декораторы ролевого доступа (app/auth/decorators.py)}, label={lst:decorators}]
def admin_required(f):
    @wraps(f)
    def wrapped(*args, **kwargs):
        if not current_user.is_authenticated or not current_user.is_admin():
            abort(403)
        return f(*args, **kwargs)
    return wrapped

def teacher_required(f):
    @wraps(f)
    def wrapped(*args, **kwargs):
        if not current_user.is_authenticated or not (
            current_user.is_teacher() or current_user.is_admin()
        ):
            abort(403)
        return f(*args, **kwargs)
    return wrapped
\end{lstlisting}

Декораторы \texttt{@admin\_required} и \texttt{@teacher\_required}
применяются к маршрутам блупринтов \texttt{admin} и \texttt{teacher}
соответственно; при нарушении прав Flask возвращает HTTP-код~403.

\subsection{Каталог курсов и расписание занятий}

Каталог (\texttt{/courses}) поддерживает одновременную фильтрацию
по языку, уровню владения и полнотекстовый поиск по названию и описанию
курса через оператор \texttt{ilike} SQLAlchemy~\cite{sqlalchemy_docs}.
Все три параметра передаются как query-строка и могут комбинироваться.

Расписание (\texttt{/schedule}) реализует недельный вид: по умолчанию
отображается текущая неделя, кнопки навигации передают параметр
\texttt{week} в формате \texttt{YYYY-MM-DD}. Занятия выбираются из базы
данных для диапазона \texttt{[week\_start, week\_start + 7 дней)} и
группируются по дням в словаре Python. Дополнительно доступна фильтрация
по языку, уровню и конкретному преподавателю; идентификаторы занятий,
на которые студент уже записан, передаются в шаблон отдельным множеством
для корректного отображения кнопок.

\subsection{Запись на занятия и контроль вместимости}

Маршрут \texttt{POST /lessons/<id>/enroll} реализует запись студента
на занятие с проверкой вместимости и дедупликацией. Если запись для
данной пары \texttt{(student\_id, lesson\_id)} уже существует
со статусом \texttt{cancelled}, она реактивируется без создания
дублирующей строки; при исчерпании мест возвращается информативное
сообщение (листинг~\ref{lst:enroll}).

\begin{lstlisting}[caption={Маршрут записи на занятие (app/main/routes.py)}, label={lst:enroll}]
@main_bp.route('/lessons/<int:lesson_id>/enroll', methods=['POST'])
def enroll(lesson_id):
    lesson = db.session.get(Lesson, lesson_id)
    existing = Enrollment.query.filter_by(
        student_id=current_user.id, lesson_id=lesson_id
    ).first()

    if existing and existing.status == EnrollmentStatus.active:
        flash('You are already enrolled in this lesson.', 'info')
    elif lesson.is_full:
        flash('This lesson is full.', 'error')
    else:
        if existing:
            existing.status = EnrollmentStatus.active
        else:
            db.session.add(Enrollment(
                student_id=current_user.id, lesson_id=lesson_id,
            ))
        db.session.commit()
        flash(f'Enrolled in {lesson.course.name}.', 'success')

    return redirect(url_for('main.schedule', week=week))
\end{lstlisting}

Свойство \texttt{lesson.is\_full} вычисляется динамически на основе
количества активных записей и поля \texttt{capacity}, не требуя
хранения отдельного счётчика в базе данных.

\subsection{Учебные материалы и экспорт успеваемости}

Учебные материалы прикрепляются к курсам с возможностью файлового
вложения. При загрузке файл сохраняется под уникальным UUID-именем
в директорию \texttt{/app/uploads/materials/}; оригинальное имя файла
фиксируется в поле \texttt{file\_name} для использования при скачивании
(листинг~\ref{lst:upload}).

\begin{lstlisting}[caption={Сохранение загружаемого файла (app/main/routes.py)}, label={lst:upload}]
UPLOAD_FOLDER = '/app/uploads/materials'

def _save_file(f):
    os.makedirs(UPLOAD_FOLDER, exist_ok=True)
    ext  = f.filename.rsplit('.', 1)[1].lower()
    path = os.path.join(UPLOAD_FOLDER, f'{uuid.uuid4().hex}.{ext}')
    f.save(path)
    return f.filename, path
\end{lstlisting}

UUID-именование файлов на диске исключает коллизии при загрузке
одноимённых файлов разными пользователями и скрывает внутреннюю
структуру хранилища.

Маршрут GET /teacher/lessons/<id>/grades/export формирует
CSV-ответ с оценками и посещаемостью всех записавшихся студентов.
Ответ отдаётся через \texttt{flask.Response} с заголовком
\texttt{Content-Disposition: attachment}, что инициирует скачивание
файла в браузере без промежуточного сохранения на сервере.

\subsection{Система уведомлений}

Внутренние уведомления хранятся в таблице \texttt{notifications}.
Сервер записывает строки при наступлении событий: одобрении или
отклонении заявки на роль преподавателя, выставлении оценки и других
действиях. Контекстный процессор передаёт счётчик непрочитанных
уведомлений во все шаблоны; он отображается на иконке колокольчика
в навигации.

На клиенте реализован JS polling: каждые 30 секунд браузер
обращается к эндпоинту \texttt{GET /notifications/api/poll}, который
возвращает JSON с общим числом непрочитанных и списком новых уведомлений
с момента последнего запроса (листинг~\ref{lst:poll}).

\begin{lstlisting}[caption={Эндпоинт polling уведомлений (app/notifications/routes.py)}, label={lst:poll}]
@notifications_bp.route('/api/poll')
@login_required
def poll():
    last_id   = session.get('last_notif_id', 0)
    new_notes = Notification.query.filter(
        Notification.user_id == current_user.id,
        Notification.is_read == False,
        Notification.id > last_id,
    ).order_by(Notification.id.asc()).all()

    if new_notes:
        session['last_notif_id'] = new_notes[-1].id

    unread_total = Notification.query.filter_by(
        user_id=current_user.id, is_read=False
    ).count()

    return jsonify({
        'unread': unread_total,
        'new': [{'id': n.id, 'message': n.message,
                 'link': n.link} for n in new_notes],
    })
\end{lstlisting}

Новые уведомления отображаются в виде всплывающих toast-сообщений
с автоматическим закрытием через 5 секунд и анимированной полосой
прогресса. Счётчик на иконке колокольчика обновляется без перезагрузки
страницы.

\subsection{Доступность исходного кода}

Исходный код веб-приложения опубликован в открытом репозитории на GitHub
и доступен по адресу:

\begin{center}
    \url{https://github.com/Miraz11287/foreign-language-center-webapp}.
\end{center}

Репозиторий содержит Flask-приложение, конфигурацию Docker Compose,
Alembic-миграции и LaTeX-документацию. Файлы загружаемых учебных
материалов исключены из репозитория через \texttt{.gitignore},
поскольку являются пользовательским контентом.

Развёрнутая версия приложения доступна по адресу:

\begin{center}
    \url{https://langcenter.digitaldrugs.tech/}.
\end{center}
